\documentclass[10pt,a4paper]{amsart}
\usepackage[margin=19mm]{geometry}
\usepackage{amsmath,amssymb,mathtools}
\usepackage[scr=rsfs,cal=euler]{mathalfa}
\usepackage{xfrac}
\usepackage[unicode,hidelinks,bookmarks=false]{hyperref}
\newcommand{\ie}{i.e.\ }
\newcommand{\cf}{cf.\ }
\newcommand{\resp}{respectively\ }
\newcommand{\Subsection}{\subsection}
\providecommand{\nobreakdash}{\nobreak\hbox{-}}
\setlength{\emergencystretch}{2em}
\multlinegap0pt
\usepackage{bm}
\usepackage{bbm}
\usepackage{dsfont}
\usepackage{xspace}
\usepackage{cancel}
\usepackage{mathtools}
\usepackage{amsthm}
\makeatletter
\@ifundefined{lemma}{\newtheorem{lemma}{Lemma}}{}
\makeatother
\newcommand{\R}{\mathbb{R}}
\title[Complete Resolution of the Butler-Costello-Graham Conjecture]{Complete Resolution of the Butler-Costello-Graham Conjecture on Monochromatic Constellations}
\author{Gang Yang, Yaping Mao}
\date{Source: arXiv:2604.02115v3 (CC BY 4.0)}
\begin{document}
\maketitle

\noindent\emph{Source and licence.} Complete Resolution of the Butler-Costello-Graham Conjecture on Monochromatic Constellations. Version arXiv:2604.02115v3 (CC BY 4.0), distributed under the Creative Commons Attribution 4.0 International licence, as recorded in the arXiv licence metadata for this exact version and shown on the abstract page https://arxiv.org/abs/2604.02115v3. Full rights notice: licenses/arxiv-2604.02115-CC-BY-4.0.txt.\par\smallskip

\noindent\textit{Editorial scope.} Counted unit: the Lemma whose source label is \texttt{lem:Riemann} in the article (source0.tex lines 508--522, source label \texttt{lem:Riemann}), delivered together with the article's own complete original proof of it (source0.tex lines 524--640, one \texttt{proof} environment of 117 lines). No mathematics is written, repaired or re-derived: statement and proof are the article's own text line for line. Required context retained uncounted: none. Every definition, display and label of those lines is retained unchanged. Omitted from this package and not redistributed: 1--507, 523--523, 641--1906. Nothing inside a retained excerpt is omitted or altered: every retained body is delivered exactly as the article's lines are written, so no omission receipt of any kind accompanies this package. Citations: the retained excerpts carry no citation command at all, so no bibliography entry of the article is transcribed and no reference is redistributed. Reference adaptation: none was needed and none was made; every reference inside the delivered unit resolves inside it, so no unresolved reference or citation is produced. Printed slips kept literal: no printed word, symbol, hypothesis or number of the counted statement or of its proof was altered. Layout and class: the article's own class is \texttt{article} with packages including amsmath, amssymb, amsthm, mathrsfs, enumitem, xcolor, geometry, hyperref; the wrapper is the lane's pinned \texttt{amsart} editorial wrapper at 10pt on A4, which restates the theorem-like environments the retained text needs and adds the article's own macro definitions that the retained text needs (\texttt{\textbackslash R}), with extra wrapper packages (bm, bbm, dsfont, xspace, cancel, mathtools). The delivered numbering is the unit's own. Assets: no retained span contains a figure environment, a graphics-inclusion command, a PDF-page inclusion, a table or a TikZ picture; no image file is copied into this package, the package declares no asset dependency and no image file is redistributed with it. Licence: version arXiv:2604.02115v3 of this article is distributed under the Creative Commons Attribution 4.0 International licence, as recorded in the arXiv licence metadata for this exact version and shown on the abstract page https://arxiv.org/abs/2604.02115v3. A full rights notice is delivered at licenses/arxiv-2604.02115-CC-BY-4.0.txt. Characters: every retained excerpt is pure ASCII, so the delivered engine, XeLaTeX with its default Latin Modern text font, reports no missing character.
\begin{lemma}\label{lem:Riemann}
Let \(f:[0,1]\to\R\) be Lipschitz.  For each
\(r\in\{0,1,\dots,D-1\}\), define \(T_r(f)=\sum_{y\in Y_r}f(y)\).  Then,
uniformly in \(r\),
\begin{equation}\label{eq:Riemann-main}
T_r(f)=\frac{n}{D}\int_0^1 f(t)\,dt+O_f(1).
\end{equation}
More precisely, if \(L_f\) is a Lipschitz constant for \(f\) and
\(M_f=\|f\|_\infty\), then
\begin{equation}\label{eq:Riemann-precise}
\left|\int_0^1 f(t)\,dt-\frac{D}{n}T_r(f)\right|
\le (L_f+M_f)\frac{D}{n}
\end{equation}
for every \(r\).
\end{lemma}

\begin{proof}
Set \(h=D/n\), and let \(N=|Y_r|\).  If \(N\ge1\), write the elements of
\(Y_r\) in increasing order as \(y_0<y_1<\cdots<y_{N-1}\).

We first treat the cases \(N=0\) and \(N=1\).  If \(N=0\), then
\(T_r(f)=0\).  This can occur only when \(n<D\), hence \(h>1\).  Therefore
\[
\left|\int_0^1 f(t)\,dt-hT_r(f)\right|
=
\left|\int_0^1 f(t)\,dt\right|
\le M_f
\le (L_f+M_f)h.
\]
If \(N=1\), say \(Y_r=\{y_0\}\), then
\[
\begin{aligned}
\left|\int_0^1 f(t)\,dt-hf(y_0)\right|
&\le
\int_0^1 |f(t)-f(y_0)|\,dt
+
|1-h|\,|f(y_0)|  \\
&\le
L_f\int_0^1 |t-y_0|\,dt
+
M_f|1-h|.
\end{aligned}
\]
Since \(N=1\), necessarily \(n<2D\), and hence \(h=D/n>1/2\).  Also
\(\int_0^1 |t-y_0|\,dt\le 1/2\le h\) and \(|1-h|\le h\).  Hence
\[
\left|\int_0^1 f(t)\,dt-hT_r(f)\right|
\le (L_f+M_f)h.
\]

It remains to consider the case \(N\ge2\).  Since consecutive integers in a
fixed residue class modulo \(D\) differ by \(D\), consecutive sampling points
satisfy
\begin{equation}\label{eq:grid-spacing}
y_{j+1}-y_j=h
\qquad (0\le j\le N-2).
\end{equation}
We partition \([0,1]\) by Voronoi cells centered at the sampling points.
Define
\[
m_0=0,
\qquad
m_j=\frac{y_{j-1}+y_j}{2}\quad (1\le j\le N-1),
\qquad
m_N=1,
\]
and set \(I_j=[m_j,m_{j+1}]\) for \(0\le j\le N-1\).  Then
\([0,1]=\bigcup_{j=0}^{N-1}I_j\).

For every \(t\in I_j\), we have \(|t-y_j|\le h\).  Indeed, for
\(1\le j\le N-2\), \eqref{eq:grid-spacing} gives
\(I_j=[y_j-h/2,y_j+h/2]\).  For the left boundary cell,
\[
I_0=[0,(y_0+y_1)/2]\subseteq [y_0-h,y_0+h],
\]
because \(0<y_0\le h\) and \(y_1=y_0+h\).  The right boundary cell is
analogous, using \(0\le 1-y_{N-1}<h\).

Moreover, every interior cell has length \(h\), and only the two boundary
cells can have length different from \(h\).  Since
\[
|I_0|=y_0+\frac h2,
\qquad
|I_{N-1}|=1-y_{N-1}+\frac h2,
\]
we have
\[
\bigl||I_0|-h\bigr|
\le \frac h2,
\qquad
\bigl||I_{N-1}|-h\bigr|
\le \frac h2.
\]
Thus
\begin{equation}\label{eq:length-deviation}
\sum_{j=0}^{N-1}\bigl||I_j|-h\bigr|\le h.
\end{equation}

By Lipschitz continuity,
\[
\left|\int_{I_j}f(t)\,dt-|I_j|f(y_j)\right|
\le
\int_{I_j}|f(t)-f(y_j)|\,dt
\le
L_fh|I_j|.
\]
Summing over \(j\), we get
\begin{equation}\label{eq:int-vs-weighted}
\left|
\int_0^1 f(t)\,dt
-
\sum_{j=0}^{N-1}|I_j|f(y_j)
\right|
\le L_fh.
\end{equation}
Next, by \eqref{eq:length-deviation} and \(|f(y_j)|\le M_f\),
\begin{equation}\label{eq:weighted-vs-riemann}
\left|
\sum_{j=0}^{N-1}|I_j|f(y_j)
-
h\sum_{j=0}^{N-1}f(y_j)
\right|
\le M_fh.
\end{equation}
Combining \eqref{eq:int-vs-weighted} and \eqref{eq:weighted-vs-riemann}, we
obtain
\[
\left|\int_0^1 f(t)\,dt-hT_r(f)\right|
\le (L_f+M_f)h.
\]
Since \(h=D/n\), this is \eqref{eq:Riemann-precise}.  Multiplying by
\(h^{-1}=n/D\) gives \eqref{eq:Riemann-main}.
\end{proof}

\end{document}
